\documentclass[11pt]{article}
\usepackage[margin=1in]{geometry}
\usepackage{amsmath,amssymb,amsthm,hyperref}
\newtheorem{proposition}{Proposition}
\setlength{\emergencystretch}{3em}
\title{A recurring obstruction to canonical endpoint projection}
\author{Research note, 2 October 2026}
\date{}
\begin{document}
\maketitle

\noindent\textbf{Scope.}
This note concerns the particular endpoint projection in
\path{FixedPartitionProjection.tex}, when its exact canonical shape is
required at every primary-switch checkpoint. It does not lower-bound
the optimal common-forest comparator, standard multi-splay's competitive
ratio, or every possible projection. The construction was proposed by
the multi-splay theory audit and independently checked against the
node-storing reference and the exact guard split word.

Let the primary reference be the perfect node BST on
$\{1,\ldots,2^h-1\}$, with $h\ge4$. Set
\[
 r=2^{h-1},\qquad y=2^{h-2},\qquad
 q=y-1,\qquad z=y+1.
\]
Thus $r$ is the reference root, $y$ its left child, $q$ the maximum
leaf of $y$'s left subtree, and $z$ the minimum leaf of $y$'s right
subtree. Let $Q$ be the all-right reference path from $y$'s left child
to $q$. Let
\[
 R=(a_1>a_2>\cdots>a_k=z),\qquad k=h-2,
\]
be the all-left reference path from $y$'s right child to $z$.
The proper-ancestor preferences on $Q$ and $R$ point toward their
respective endpoints, and $r$ prefers left. Only $y$'s preference
changes when the requests alternate between $q$ and $z$.
The final switches at the requested reference leaves change no
nonempty auxiliary membership; their external directions can be
included as a two-period harmless bit state.

Choose a global comparator BST in which $q$ is the root and $z$ is
its right child. Key $y$, the only key strictly between $q$ and $z$,
is the left child of $z$. Put
\[
 a_{k-1}<a_{k-2}<\cdots<a_1<r
\]
on an increasing right spine below $z$. All remaining actual keys
can be placed in the opaque inorder gaps of this displayed skeleton;
in particular the keys of $Q$ other than $q$ are below $q$'s left
subtree. Root rotations between $q$ and $z$ serve the alternating
sequence with constant cost per access and preserve the displayed
cold right spine. No cold key is requested.

\begin{proposition}
Requiring the exact canonical endpoint-projection shape for this
global comparator and the current preferred partition makes the
prescribed common-forest switch word cost $\Omega(h)$ per two original
requests, although the global comparator cost and the number of
primary switches are both $O(1)$ per two requests.
\end{proposition}

\begin{proof}
Immediately before a request to $q$, the root primary component has
actual key set
\[
 C_R=\{r,y\}\cup R.
\]
The incoming left reference subtree of $y$ is one exit block of this
component. Its two selected external endpoints have global lowest
common ancestor $q$, so their ghost node is the global root. In the
conversion to actual component separators, this leftmost exit block
has only the upper neighbor $y$. The projected component therefore
has $y$ at its root. Its right subtree contains $z$ followed by the
displayed right spine through $a_{k-1},\ldots,a_1,r$.

These are real component separators, not removable ghost nodes. There
are no selected interior exit-block representatives between consecutive
nodes of this spine that could suppress a displayed separator: all
displayed actual keys belong to $C_R$. The optional ghost images in
their inorder gaps preserve the displayed root-to-$r$ ancestor chain.
In particular, $r$ has depth at least $k$ in the right daughter of $y$.

The standard right-to-left switch at $y$ has no lower ancestor guard,
and has upper ancestor guard $r$. Its exact forest word begins with
$\operatorname{Split}(y)$ and then $\operatorname{Split}(r)$ in the
right daughter. The first pivot is already exposed, but the latter
requires a root-to-$r$ traversal and hence $\Omega(k)$ pointer cost.
Preparing $r$ before the first split cannot avoid this cost: its
initial depth is $\Omega(k)$, and a rooted execution must reach it.

After the switch the root component is
$C_Q=\{r,y\}\cup Q$, while $R$ is parked. In its canonical projection,
the entire old right reference subtree is now one exit block. Its
endpoint ghost is rooted at $z$, which is at depth one in the chosen
global comparator. The conversion exposes its two boundary keys
$y,r$ at constant depth. Thus the coarser projection puts $r$ shallow.

On the next request to $z$, one global root rotation again suffices.
The primary switch attaches $R$ again. Requiring the exact canonical
projection of $C_R$ restores the long real-separator chain ending at
$r$. If this restoration itself costs linear work, that work is
already charged to the same two-request period; otherwise the next
$q$ request encounters the linear-depth guard just proved. Either
way every period has $\Omega(k)=\Omega(h)$ cost.
\end{proof}

The global tree, preferred membership partition, and canonical
projected shapes return after two original requests. If the analytic
external leaf directions are also tracked, use four requests. Hence
no potential depending only on this complete canonical checkpoint
state can amortize its recurrent linear seam cost against the constant
source work and constant switch count. A one-time initialization term
does not pay arbitrarily many repetitions.

An unrestricted forest comparator can instead retain $r$ near the root
when $R$ is attached again. This is legal and avoids restoring the bad
canonical shape. Likewise one could choose a different, guard-aware
global comparator for the same request sequence. These possibilities
are why this proposition refutes the specific canonical transport,
not a bound of the form
$C_{\mathrm{forest}}(W_{\mathrm{MS}})=O(\mathrm{OPT}+n)$.

\end{document}
